% SOURCE_AUTHORITY: sga5_fr_workpass.tex, lines 9477--9520 (LF-normalized unit).
% SOURCE_SHA256: 791F4EFFC5E02832D5D77ED03518C8156D6F07E4C8238B03545DB93D883FBB28
\emph{Demostración de 3.4.} (según A. Grothendieck: \emph{La théorie des classes de Chern}. Bull. S.M.F. (1958), pp. 137--154).

(i) Sea $E$ un $\mathcal O_S$-módulo localmente libre; se pone $P=P_S(\check E)$ y se denota por $p:P\to S$ la proyección canónica. Recordemos que existe un homomorfismo canónico sobreyectivo
\[
p^*(\check E)\longrightarrow \mathcal O_P(1).
\]
En particular, si $E$ es invertible, poniendo $E=L$ de acuerdo con las notaciones de (i), se tiene $P=S$ y un isomorfismo
\[
L\xrightarrow{\sim}\mathcal O_P(1).
\]
Por tanto, $\xi+[L]=0$ y la afirmación se sigue de la definición de las clases de Chern.

(ii) Verificación inmediata que se deja al lector.

(iii) Utilizando una partición de $S$ en abiertos sobre los cuales los módulos considerados tienen rango constante, se puede suponer que $E,E',E''$ tienen rango constante igual a $r,r',r''$, respectivamente. Supongamos entonces demostrado el lema siguiente y veamos cómo implica la afirmación.

\begin{statement}{Lema 3.6}
Sea $E$ un $\mathcal O_S$-módulo localmente libre dotado de una filtración finita
\[
0=E_0\subset E_1\subset\cdots\subset E_i\subset\cdots\subset E_r=E
\]
cuyos cocientes consecutivos $L_i=E_i/E_{i-1}$ $(1\le i\le r)$ son invertibles. Entonces se verifica la igualdad siguiente:
\[
c(E)=(1+[L_1])(1+[L_2])\cdots(1+[L_r]).
\]
\end{statement}

Sea $h:Z\to S$ el producto fibrado de los esquemas de banderas de $E'$ y $E''$. Por la definición de los esquemas de banderas, es claro que $h^*(E')$ y $h^*(E'')$ admiten filtraciones finitas con cocientes invertibles $(L'_i)$ $(1\le i\le r')$ y $(L''_j)$ $(1\le j\le r'')$, respectivamente. Se sigue que $h^*(E)$ admite una filtración cuyos cocientes consecutivos son los $L'_i$ y los $L''_j$. La aplicación del lema 3.6 muestra entonces que
\[
c(h^*(E))=
\prod_{1\le i\le r'}c(L'_i)\prod_{1\le j\le r''}c(L''_j)
=c(h^*(E'))c(h^*(E'')),
\]
es decir, gracias a (ii),
\[
h^*(c(E))=h^*(c(E')c(E'')).
\tag{3.5.1}
\]
Pero el morfismo $h$ es una composición de morfismos del tipo
\[
P_T(F)\longrightarrow T
\qquad (F\text{ }\mathcal O_T\text{-módulo localmente libre}),
\]
por lo que, según (2.2.5), el homomorfismo de imagen inversa $h^*:A^*(S)\to A^*(Z)$ es inyectivo. De aquí se sigue la afirmación deseada gracias a (3.5.1).
